% Pista geo-23j2-q3 · Geometria
% Procedència: PAU Catalunya 2023 · Convocatòria de juny · Sèrie 5 · Qüestió 3
\documentclass[11pt,a4paper]{article}
\usepackage{pista}
\pistainfo{Catalunya 2023}{Convocatòria de juny}{Sèrie 5}

\begin{document}

\section*{\textcolor{blauPista}{Qüestió 3}}

\noindent Considereu les rectes a l'espai $r: x = -y = z + m$ i $s: \begin{cases} x + y = 1 \\ x - z = 0 \end{cases}$, en què $m$ és un paràmetre real.

\subsection*{\textit{a)} \normalfont Estudieu la posició relativa per als diferents valors del paràmetre $m$. \hfill\textnormal{[1,25~punts]}}

\pista{Treu de cada recta un vector director i un punt. És bastant senzill obtenir l'equació contínua de $r$, però vigila, perquè surten alguns signes negatius. De $s$, com que ve donada com a intersecció de dos plans, el vector director és el producte vectorial dels dos vectors normals. Compara els dos vectors directors: si són proporcionals, les rectes només poden ser paral·leles o coincidents, i ho decideixes mirant si un punt d'una pertany a l'altra. Comprova si la condició de coincidència es pot complir per algun valor de $m$: mira si totes les igualtats es poden complir alhora.}

\subsection*{\textit{b)} \normalfont Calculeu el valor de $m$ perquè la distància entre les rectes $r$ i $s$ sigui de $\sqrt{2}$ unitats. \hfill\textnormal{[1,25~punts]}}

\pista{Com que són paral·leles, fes servir la fórmula de la distància $d(r, s) = \dfrac{|\vec{v} \times \vec{PQ}|}{|\vec{v}|}$, amb $P \in r$ i $Q \in s$. Et quedarà una expressió que depèn de $m$; iguala-la a $\sqrt{2}$, eleva al quadrat per a treure l'arrel i resol l'equació de segon grau resultant en $m$ (sortiran dos valors).}

\end{document}
